\chapter{Conclusão}

Este trabalho apresentou o Spectrum, uma plataforma \textit{web} moderna e extensível para a gestão de modelos de regressão simbólica. Através da implementação de uma arquitetura baseada em \textit{micro-packages} e do uso de tecnologias de ponta como FastAPI, React e RabbitMQ, foi possível criar um ambiente robusto que supera limitações de ferramentas tradicionais, oferecendo execução remota e gerenciamento de \textit{pipelines}.

A principal contribuição inovadora deste projeto reside na introdução da \textit{Inference Query Language} (IQL), que preenche a lacuna entre a busca automatizada de expressões e o conhecimento analítico do pesquisador. A capacidade de consultar o espaço de soluções de forma declarativa transforma a regressão simbólica de um processo puramente computacional em uma ferramenta de descoberta científica assistida.

Como trabalhos futuros, propõe-se a expansão do suporte a outros \textit{backends} de regressão simbólica (como o PySR), a melhoria das capacidades de visualização de dados no \textit{front-end} e o refinamento da gramática do IQL para suportar otimizações automáticas baseadas em padrões matemáticos. Com essas adições, o Spectrum tem o potencial de se tornar uma ferramenta padrão para cientistas que buscam equilíbrio entre precisão e interpretabilidade.
